\documentclass[11pt]{article}
\usepackage[T1]{fontenc}
\usepackage{lmodern}
\usepackage[margin=1.05in]{geometry}
\usepackage{amsmath,amssymb,amsthm,mathtools}
\usepackage{microtype}
\usepackage{booktabs}
\usepackage{enumitem}
\usepackage{hyperref}
\usepackage{xurl}
\hypersetup{hidelinks,pdftitle={An Exact Uniform Component and Analytic Mod-Poisson Convergence for Luce Permutations},pdfauthor={Christopher D. Long},pdfsubject={Draft research manuscript}}
\numberwithin{equation}{section}
\newtheorem{theorem}{Theorem}[section]
\newtheorem{proposition}[theorem]{Proposition}
\newtheorem{lemma}[theorem]{Lemma}
\newtheorem{corollary}[theorem]{Corollary}
\newtheorem{inputtheorem}[theorem]{Input theorem}
\theoremstyle{definition}
\newtheorem{definition}[theorem]{Definition}
\theoremstyle{remark}
\newtheorem{remark}[theorem]{Remark}
\newcommand{\E}{\mathbb E}
\newcommand{\Pp}{\mathbb P}
\newcommand{\one}{\mathbf 1}
\newcommand{\dd}{\,\mathrm d}
\newcommand{\TV}{\mathrm{TV}}
\newcommand{\Exp}{\operatorname{Exp}}
\newcommand{\Pois}{\operatorname{Pois}}
\newcommand{\Unif}{\operatorname{Unif}}
\newcommand{\Var}{\operatorname{Var}}
\newcommand{\Tr}{\operatorname{Tr}}
\newcommand{\supp}{\operatorname{supp}}
\newcommand{\cF}{\mathcal F}
\newcommand{\cH}{\mathcal H}
\newcommand{\cG}{\mathcal G}
\newcommand{\cA}{\mathcal A}
\newcommand{\cB}{\mathcal B}
\newcommand{\Kop}{\mathsf K}
\newcommand{\Qop}{\mathsf Q}
\newcommand{\Top}{\mathsf T}
\newcommand{\eps}{\varepsilon}
\newcommand{\pdual}{{p^{\prime}}}
\newcommand{\mres}{\mathfrak m}
\newcommand{\cyc}{\operatorname{cyc}}
\newcommand{\norm}[1]{\left\|#1\right\|}
\newcommand{\abs}[1]{\left|#1\right|}
\title{An Exact Uniform Component and Analytic\\Mod-Poisson Convergence for Luce Permutations}
\author{Christopher D. Long\\\texttt{galizur@gmail.com}}
\date{Draft --- 20 October 2026}

\begin{document}
\maketitle
\begin{abstract}
Let $K_n$ be the total number of cycles of a Luce permutation generated by independent exponential clocks. For rates with a uniformly bounded ratio, we prove a locally uniform complex bound for $n^{1-z}\E z^{K_n}$ on a disk about $1$, with radius depending only on the comparability constant. For rates converging uniformly to a positive continuous profile $f$, this yields
\[
 n^{1-z}\E z^{K_n}\longrightarrow
 \frac{\exp\{A_f(z-1)\}}{\Gamma(z)},
 \qquad
 A_f=-\log\det\bigl(I-(\Kop-\Pi)\bigr),
\]
locally uniformly near $1$, where $\Kop$ is the microscopic Luce operator and $\Pi$ projects onto constants. The proof extracts an exactly uniform conditional permutation by common-rate thinning of the clocks. A real cycle-fugacity estimate supplies an exponential moment for the residual cycle count; the exact uniform factor then retains the complex cancellation. The residual cycles converge to a Poisson law, whose parameter is identified by a rank-one Fredholm identity. The proof uses the previously established microscopic local-entry and envelope results and the unweighted uniform-root macroscopic limit. It does not use higher-cumulant peeling, finite-deletion independence, or complex-tilted endpoint estimates.
\end{abstract}

\setcounter{tocdepth}{1}
\tableofcontents
\newpage
\raggedbottom

\section{Model, statements, and proof dependencies}\label{sec:introduction}

The cycle structure of the Luce model was raised as an open problem by
Chatterjee--Diaconis--Kim \cite[Section~1]{CDK}. Borga--Chatterjee--Diaconis
\cite[Section~5.2]{BCD} formulate a fixed-point Poisson conjecture and explicitly
ask for the joint cycle counts and the total number of cycles. The same broad
question was highlighted by Diaconis in his Stanford combinatorics seminar
\cite{DiaconisSeminar}. We answer the total-cycle question for rates converging
uniformly to a positive continuous profile, in the stronger form of locally
uniform analytic mod-Poisson convergence. The correction to the uniform law is
a profile-dependent Fredholm determinant; the limiting higher factorial
cumulants are universal. We do not attribute this particular analytic
mod-Poisson formulation to those authors.

The result complements the fixed short-cycle Poisson limits of \cite{LongMicro}
and the macroscopic $\mathrm{PD}(1)$ law of \cite{LongMacro}, under their respective
hypotheses. Earlier total-cycle Poisson and lattice Edgeworth estimates
\cite{LongTotal} left the profile-dependent mean constant and analytic
mod-Poisson convergence separate. The present argument identifies the analytic
residue; it does not replace the quantitative lattice error estimates in that
paper. Its main structural ingredient is an exactly uniform conditional
permutation extracted by thinning the exponential clocks.

Let $X_{n,i}\sim\Exp(\theta_{n,i})$, $1\le i\le n$, be independent, and define the label-to-rank permutation
\[
 \sigma_n(i)=1+\#\{j:X_{n,j}<X_{n,i}\}.
\]
Write $C_{n,r}$ for its number of $r$-cycles, $L_n(i)$ for the length of the cycle containing $i$, and
\[
 K_n=\sum_{r=1}^n C_{n,r},\qquad
 G_n(z)=\E z^{K_n},\qquad H_n=\sum_{j=1}^n\frac1j.
\]
The inverse permutation $\pi_n=\sigma_n^{-1}$ is the draw order obtained by sampling labels without replacement, with probabilities proportional to their remaining rates. Inversion preserves the cycle counts. This distinction between rank and draw order is important in the sequential argument below; see Luce \cite{Luce}, Chatterjee--Diaconis--Kim \cite{CDK}, and Borga--Chatterjee--Diaconis \cite{BCD}.

A common scaling of the rates does not affect the permutation. For the uniform estimates we assume
\begin{equation}\label{eq:comparable}
 1\le\theta_{n,i}\le\varkappa<\infty.
\end{equation}
For the limiting residue we additionally assume, after such a scaling, that
\begin{equation}\label{eq:profile}
 \max_{1\le i\le n}|\theta_{n,i}-f(i/n)|\longrightarrow0,
 \qquad f\in C[0,1],\quad 1\le f\le\varkappa.
\end{equation}
The usual assumption $\theta_{n,i}=a_n(f(i/n)+o(1))$ with $0<\min f\le\max f<\infty$ reduces to this form, with a fixed enlarged $\varkappa$, after ignoring finitely many $n$.

Set
\begin{align}
 \cA(t)&=\int_0^1(1-e^{-f(x)t})\dd x,
 &\cB(t)&=\cA'(t)=\int_0^1 f(x)e^{-f(x)t}\dd x,\label{eq:profile-transforms}\\
 \tau&=\cA^{-1},
 &\rho(x,y)&=\frac{f(x)e^{-f(x)\tau(y)}}{\cB(\tau(y))},\qquad 0<y<1.\label{eq:rho}
\end{align}
Let $\Kop$ be the integral operator with kernel $\rho$ on $L^2[0,1]$, and let $\Pi h=\int_0^1h(x)\dd x$. Define
\begin{equation}\label{eq:lambda}
 \lambda_r=\frac1r\int_{[0,1]^r}\rho(x_1,x_2)\cdots\rho(x_r,x_1)
 \dd x_1\cdots\dd x_r,
 \qquad r\ge1,
\end{equation}
with the integral interpreted by increasing bulk truncations when necessary.

\begin{theorem}[Uniform analytic control]\label{thm:uniform}
For each $\varkappa<\infty$, there are $\delta_\varkappa>0$ and $C_\varkappa<\infty$ such that, for every array satisfying \eqref{eq:comparable},
\begin{equation}\label{eq:normal-bound}
 \sup_{n\ge1}\sup_{|z-1|\le\delta_\varkappa}
 \abs{n^{1-z}G_n(z)}\le C_\varkappa.
\end{equation}
Moreover,
\begin{equation}\label{eq:uniform-cumulants}
 |\E K_n-H_n|\le C_\varkappa,
 \qquad
 |\kappa_k^{\mathrm F}(K_n)|\le C_\varkappa^k(k-1)!,\quad k\ge2,
\end{equation}
where $\kappa_k^{\mathrm F}(K_n)$ is the $k$th derivative of $\log G_n(1+w)$ at $w=0$. The constants are independent of the arrangement of the rates and of $k$.
\end{theorem}

\begin{theorem}[Profile-dependent analytic mod-Poisson convergence]\label{thm:profile}
Under \eqref{eq:comparable}--\eqref{eq:profile}, the operator $\Kop$ is trace class, the series
\begin{equation}\label{eq:Af}
 A_f=\sum_{r\ge1}\left(\lambda_r-\frac1r\right)
     =-\log\det\bigl(I-(\Kop-\Pi)\bigr)
\end{equation}
is absolutely convergent, and
\begin{equation}\label{eq:mod-poisson}
 n^{1-z}\E z^{K_n}\longrightarrow
 \frac{e^{A_f(z-1)}}{\Gamma(z)}
\end{equation}
locally uniformly on a complex disk $|z-1|<\delta_\varkappa$. In particular,
\begin{equation}\label{eq:exact-mean-residue}
 e^{-(z-1)\E K_n}\E z^{K_n}
 \longrightarrow \frac{e^{-\gamma(z-1)}}{\Gamma(z)}
\end{equation}
locally uniformly near $1$.
\end{theorem}

Here and below the logarithm of a nonvanishing Fredholm determinant is chosen by continuation from the identity along the indicated real contraction. The limiting determinants used below are positive. The conclusion is the \emph{local analytic} form of mod-Poisson convergence: with $z=e^w$, it is a holomorphic convergence statement for $w$ in a neighborhood of $0$. We do not assert a result on the entire imaginary axis. For the mod-Poisson framework see Kowalski--Nikeghbali \cite{KN}; for cycle-weighted, conjugation-invariant permutation measures see Nikeghbali--Zeindler \cite{NZ}. The Luce law is not generally conjugation invariant.

The profile correction and the coefficientwise conjectured residue appear in the unpublished manuscript \cite{LongMean}. The present proof identifies the analytic limit without importing that manuscript's higher-factorial-cumulant theorem. The two prior model-specific inputs actually needed are stated next.

\subsection{The two imported asymptotic inputs}\label{subsec:inputs}

\begin{inputtheorem}[Microscopic local entry and summable envelopes]\label{input:micro}
Under \eqref{eq:profile}, fix $k\ge1$, $0<\eps<1/2$, and $\eta>0$.
Let $\mathcal B_{\eps,\eta,n}^{(k)}$ be the distinct target tuples
$\mathbf j=(j_1,\ldots,j_k)$ satisfying
\[
 j_a/n\in[\eps,1-\eps],\qquad |j_a-j_b|\ge\eta n\quad(a\ne b).
\]
Uniformly over $\mathbf j\in\mathcal B_{\eps,\eta,n}^{(k)}$ and all distinct
source labels $\mathbf i=(i_1,\ldots,i_k)$ in $[n]$, with no bulk or separation
condition on the sources,
\begin{equation}\label{eq:local-input}
 \Pp\{\sigma_n(i_a)=j_a,\ 1\le a\le k\}
 =n^{-k}\prod_{a=1}^k\rho(i_a/n,j_a/n)+o(n^{-k}).
\end{equation}
Moreover, the summable envelopes have the supremum-over-sources form
\begin{equation}\label{eq:envelope-input}
 \lim_{\eps\downarrow0}\lim_{\eta\downarrow0}\limsup_{n\to\infty}
 \sum_{\substack{\mathbf j\in[n]^k\ \mathrm{distinct}\\
                 \mathbf j\notin\mathcal B_{\eps,\eta,n}^{(k)}}}
 \ \sup_{\mathbf i\in[n]^k\ \mathrm{distinct}}
 \Pp\{\sigma_n(i_a)=j_a,\ 1\le a\le k\}=0.
\end{equation}
Consequently, the fixed cycle counts have joint independent Poisson limits
with means \eqref{eq:lambda}, and all fixed mixed factorial moments converge.
\end{inputtheorem}

The local statement is the proposition ``Separated bulk local estimate'' in
\cite[Section~3]{LongMicro}. The lemmas ``Endpoint domination'' and
``Bulk collision domination'', together with the corollary ``Endpoint tightness
and improper integrability'', give \eqref{eq:envelope-input}. The order of limits
is first $n\to\infty$ with fixed $\eps,\eta$, then $\eta\downarrow0$, and finally
$\eps\downarrow0$. For a prescribed cycle collection the sources and targets
are the same vertex set, cyclically reordered. Its omitted contribution is
therefore bounded, up to a fixed combinatorial factor, by the displayed
sum over targets of the source supremum. Since adding thinning marks only
decreases each summand, this same envelope applies to the marked cycles in
Section~\ref{sec:marked}.

\begin{inputtheorem}[Unweighted macroscopic root law]\label{input:root}
For every array satisfying \eqref{eq:comparable} and an independent uniform label $I_n$, under the original, untilted measure,
\begin{equation}\label{eq:root-input}
 L_n(I_n)/n\ \Rightarrow\ \Unif(0,1).
\end{equation}
\end{inputtheorem}

This is the uniform-root consequence of the macroscopic cycle universality theorem in \cite[Theorem 2.1]{LongMacro}. It follows there from the $\mathrm{PD}(1)$ limit; see also \cite[Chapter 3]{Pitman} for the standard size-biased interpretation. Only \eqref{eq:root-input} is used here. No residual-history or finite-deletion version is imported.

\paragraph{Notation.}
The rate-comparability bound is $\varkappa$, whereas
$\kappa_k^{\mathrm F}(Y)$ denotes a factorial cumulant of order $k$.
The conjugate H\"older exponent is $\pdual=p/(p-1)$.
In the endpoint argument, $\Delta=r/n$, $\mres=\sqrt{nr}$,
$t_{\mathrm{cut}}$ is a clock-time cutoff, and $\omega$ is the tail fraction.
The limiting kernels are $\rho$ and $\rho_{\mathrm{res}}$; the latter is
introduced with the residual cycles. The symbol $\delta_{\varkappa}$ is
reserved for the analytic radius.

\subsection{Proof organization}
Theorem~\ref{thm:uniform} is proved independently of Inputs~\ref{input:micro}--\ref{input:root}. Sections~\ref{sec:switch}--\ref{sec:real} establish an optimal upper bound for $G_n(s)$ at one fixed real $s>1$. Section~\ref{sec:uniform-core} extracts an exact uniform conditional permutation and obtains the complex bound. Under the two asymptotic inputs, Sections~\ref{sec:operator}--\ref{sec:marked} identify the remaining finite-intensity cycle cloud and finish Theorem~\ref{thm:profile}.

\section{Exact switching and real-fugacity integrability}\label{sec:switch}

Throughout the uniform part, constants depend only on $\varkappa$ unless another dependence is displayed. We suppress $n$ on the clocks and rates. For real $t>0$, define
\begin{equation}\label{eq:tilt}
 W_{n,t}=\frac{t^{K_n}}{G_n(t)},\qquad
 \dd\Pp_t=W_{n,t}\dd\Pp,
\end{equation}
and write $\E_t$ for tilted expectation.

\subsection{A sequential upper bound}
\begin{lemma}[Crude exponential moments]\label{lem:crude-mgf}
For $u\ge1$,
\begin{equation}\label{eq:crude-mgf}
 G_n(u)\le\prod_{m=1}^n\left(1+\frac{\varkappa(u-1)}m\right)
 \le e^{\varkappa(u-1)H_n}.
\end{equation}
\end{lemma}
\begin{proof}
Expose the draw-order permutation $\pi_n=\sigma_n^{-1}$ by assigning its images in source order $1,\ldots,n$. Conditional on the first $i-1$ images, its next image is chosen from the remaining labels with probabilities proportional to their rates. The partial directed graph consists of paths and completed cycles. The current source $i$ has no assigned outgoing edge, so it is the end of one path, possibly an isolated vertex. Exactly one remaining target, the beginning of that path, closes a cycle. If $m=n-i+1$ targets remain, its conditional probability is at most $\min(1,\varkappa/m)$. Each final cycle is closed exactly once. If $B_i$ is its closure indicator, then $K_n=\sum_iB_i$ and
\[
 \E[u^{B_i}\mid\pi_n(1),\ldots,\pi_n(i-1)]
 \le1+\varkappa(u-1)/m.
\]
Iterating conditional expectations proves the product bound; $1+x\le e^x$ proves the second inequality.
\end{proof}

Fix once and for all
\begin{equation}\label{eq:exponents}
 p=1+\frac1{4\varkappa},\qquad \pdual=\frac p{p-1}=4\varkappa+1,
 \qquad s=p^{1/\pdual}>1.
\end{equation}
For $1\le t\le s$, $G_n(t)\ge1$, $t^\pdual\le p$, and Lemma~\ref{lem:crude-mgf} gives
\begin{equation}\label{eq:weight-norm}
 \norm{W_{n,t}}_\pdual\le e^{H_n/(4\pdual)}\le C_\varkappa n^{1/(4\pdual)}.
\end{equation}
We will use this only against quantities with a genuine negative power of $n$. It is not a uniform bound for the change-of-measure density.

\subsection{The master switch}
For $i,j\in[n]$, put
\begin{equation}\label{eq:R}
 R_{ij}=\exp\{-(\theta_i-\theta_j)(X_j-X_i)\},\qquad
 S_{n,r}=\sum_{i:L_n(i)>r}R_{i,\sigma_n^r(i)},\quad r<n.
\end{equation}

\begin{lemma}[Exact marked switch]\label{lem:switch}
For every function $F$ on permutations,
\begin{equation}\label{eq:master}
 r(n-r)\E[C_{n,r}F(\sigma_n)]
 =\E\sum_{i:L_n(i)>r}R_{i,\sigma_n^r(i)}
 F\bigl(\sigma_n\circ(i\ \sigma_n^r(i))\bigr).
\end{equation}
In particular,
\begin{equation}\label{eq:tilted-switch}
 r(n-r)\E_tC_{n,r}=t\E_tS_{n,r}.
\end{equation}
\end{lemma}
\begin{proof}
Swapping clocks $X_i,X_j$ sends $\sigma$ to $\sigma\circ(i\ j)$ and changes the density by $R_{ij}$. On $\{j=\sigma^r(i),L(i)>r\}$, the transposition splits one cycle into two, the cycle containing $j$ having length $r$. The reverse configuration consists of a label $j$ in an $r$-cycle and any $i$ outside that cycle. There are $r(n-r)$ reverse marks per $r$-cycle. Sum the change-of-variables identity over the ordered marks. This proves \eqref{eq:master}, as in \cite[Marked master switching identity]{LongMeso}. On every forward split, $K$ increases by one. Substitution of $F=t^{K_n}$ and division by $G_n(t)$ proves \eqref{eq:tilted-switch}. All functions are bounded for each fixed $n$; integrability of the swap factors follows also from the next lemma.
\end{proof}

\begin{lemma}[Tilted selected-factor moments]\label{lem:selected}
With $p,s$ as in \eqref{eq:exponents},
\begin{equation}\label{eq:selected-moment}
 \sup_{n\ge1,\,r<n,\,1\le t\le s}
 \frac1n\E_t\sum_{i:L_n(i)>r}R_{i,\sigma_n^r(i)}^p<\infty.
\end{equation}
The original-measure moments $\sup_{n,i,j}\E R_{ij}^p$ are finite and $\E R_{ij}=1$.
\end{lemma}
\begin{proof}
Set $A_i=e^{(\theta_i-1)X_i}$. If $\theta_i\ge\theta_j$, then $R_{ij}\le A_i$; the reverse inequality of rates gives $R_{ij}\le A_j$. Thus $R_{ij}^p\le A_i^p+A_j^p$. The map $i\mapsto\sigma^r(i)$ permutes the labels in cycles longer than $r$, whence
\begin{equation}\label{eq:deterministic-domination}
 \frac1n\sum_{i:L_n(i)>r}R_{i,\sigma_n^r(i)}^p
 \le\frac2n\sum_i A_i^p.
\end{equation}
When $\varkappa>1$, the choice of $p$ satisfies $p^2<\varkappa/(\varkappa-1)$. Therefore
\[
 \E A_i^{p^2}=\frac{\theta_i}{\theta_i-p^2(\theta_i-1)}\le C_\varkappa.
\]
For $\varkappa=1$, all $A_i=R_{ij}=1$.

Let $V_n=n^{-1}\sum_i A_i^p$. The independent-sum inequality for $1\le p\le2$ gives
\[
 \norm{V_n-\E V_n}_p\le C_\varkappa n^{-1/\pdual}.
\]
A short proof of the version needed here is to symmetrize centered independent summands with independent copies, insert independent Rademacher signs, use the conditional $L^p$--$L^2$ inequality, and then use $(\sum a_i^2)^{p/2}\le\sum|a_i|^p$. This yields $\E|\sum Z_i|^p\le2^p\sum\E|Z_i|^p$; compare von Bahr--Esseen \cite{vBE}.
By H\"older and \eqref{eq:weight-norm},
\[
 \E_t|V_n-\E V_n|\le C_\varkappa n^{-3/(4\pdual)}.
\]
Since $\sup_n\E V_n<\infty$, \eqref{eq:deterministic-domination} proves \eqref{eq:selected-moment}. The original independent-pair moment follows from the same domination. For $i\ne j$, independence and the exponential moment formula give
\[
 \E R_{ij}
 =\E e^{(\theta_i-\theta_j)X_i}
  \E e^{-(\theta_i-\theta_j)X_j}
 =\frac{\theta_i}{\theta_j}\frac{\theta_j}{\theta_i}=1.
\]
The case $i=j$ is immediate.
\end{proof}

\begin{corollary}[Tilted short-root escape]\label{cor:root-escape}
Uniformly for $1\le t\le s$,
\begin{equation}\label{eq:cycle-crude}
 \E_t C_{n,r}\le\frac{C_\varkappa n}{r(n-r)},\quad r<n,
 \qquad
 \Pp_t\{L_n(I_n)\le a\}\le C_\varkappa a/n,\quad a\le n/2,
\end{equation}
where $I_n$ is uniform and independent of the clocks under $\Pp_t$.
\end{corollary}
\begin{proof}
Lemma~\ref{lem:selected} bounds $\E_tS_{n,r}/n$ uniformly. Apply \eqref{eq:tilted-switch}, then use
$\Pp_t(L_n(I_n)=r)=r\E_tC_{n,r}/n$ and sum for $r\le a$.
\end{proof}

\section{Hereditary flatness under a cycle fugacity}\label{sec:row}

Partition the clock-time axis into intervals. Let $H_a$ be the labels in interval $a$ and $Q_a$ the corresponding consecutive ranks. The coarse sigma-field $\cH$ records all the sets $H_a$ and the ordered but unlabeled clock values in each interval. Conditional on $\cH$, a full assignment has original weight proportional to
\begin{equation}\label{eq:matching-weight}
 \prod_i e^{-\theta_i x_{\sigma(i)}},
\end{equation}
where $x_j$ is the fixed $j$th ordered clock. Conditioning on compatible revealed assignments fixes their factors. Under $\Pp_t$, the additional factor is $t^{K(\sigma)}$. We do not assume that residual rows are independent under this tilt.

\begin{lemma}[Single-row contraction and insertion]\label{lem:row-flatness}
Let $\cF$ be the sigma-field generated by $\cH$ and the compatible revealed
source--target assignments, with no additional constraint on unexposed edges.
In a finite row, let $S_a\subseteq H_a$ and $T_a\subseteq Q_a$ be the residual
source and target sets. Suppose $|S_a|=|T_a|=m\ge1$ and the row has clock-time
width at most $h\le1$. For each $x\in S_a$, uniformly for $1\le t\le s$,
\begin{align}
 d_{\TV}\bigl(\mathcal L_t(\sigma(x)\mid\cF),\Unif(T_a)\bigr)
 &\le C_\varkappa\left(h+\frac1m\right),\label{eq:row-TV}\\
 \Pp_t(\sigma(x)=y\mid\cF)&\le C_\varkappa/m,
 \qquad y\in T_a.\label{eq:row-upper}
\end{align}
A sharper version of \eqref{eq:row-TV} is $C_\varkappa h+C_\varkappa|t-1|/m$.
\end{lemma}
\begin{proof}
Fix all unexposed assignments outside this row. The fixed edges form paths from $T_a$ to $S_a$, including length-zero paths, and closed cycles disjoint from $S_a$. Following the paths gives a bijection $\tau:T_a\to S_a$. For the remaining row bijection $\pi$, put $g=\tau\circ\pi$. Then
\[
 K(\sigma)=K_{\mathrm{fixed}}+K(g),\qquad
 W(g)=\prod_{v\in S_a}e^{-\theta_v x_{\tau^{-1}(g(v))}}.
\]
Thus the conditional law of $g$ is proportional to $t^{K(g)}W(g)$.

For $m\ge2$, fix $j\ne x$. If $g(x)=x$, let $v=g^{-1}(j)$ and interchange the images of $x$ and $v$. This is a bijection onto the permutations satisfying $g(x)=j$. It merges the singleton $x$ with another cycle, reducing the cycle count by one. The logarithmic change of $W$ has absolute value at most $\varkappa h$. Therefore, with $p_j=\Pp_t(g(x)=j)$ and $\beta=\varkappa h$,
\begin{equation}\label{eq:insertion-ratio}
 e^{-\beta}\le\frac{p_x}{t p_j}\le e^\beta,\qquad j\ne x.
\end{equation}
Normalization shows that this law is within $\tfrac12(e^{2\beta}-1)$ of the probability vector with mass $t/(m-1+t)$ at $x$ and mass $1/(m-1+t)$ elsewhere. Its distance from uniform is
\[
 \frac{(m-1)|t-1|}{m(m-1+t)}.
\]
The coordinate upper bound follows from the same ratios. The case $m=1$ is immediate. Transfer by $\tau^{-1}$ and then average over the fixed outside-row assignments. Total variation is convex and the bounds are uniform in those assignments, proving the assertions.
\end{proof}

\begin{remark}\label{rem:conditioning}
The lemma applies at an adaptively selected tip because the revealed history specifies only compatible assignments. We never condition on the future event that a path will avoid closure. Such events are handled as failures using Corollary~\ref{cor:root-escape}. This is also why conditioning on outside-row completions in the proof causes no circularity: it is an auxiliary conditioning used only to obtain a bound uniform in all those completions.
\end{remark}

\section{A finite one-root endpoint estimate}\label{sec:endpoint}

We give the real-tilted version of the coarse-environment construction of \cite{LongMeso}. Only one root and a first-moment observable comparison are required. In particular, no two-root or selected second-moment estimate is used.

Fix $r<n$ and an integer terminal depth $\ell$. Put
\begin{equation}\label{eq:endpoint-parameters}
 \begin{gathered}
 \Delta=\frac rn,\qquad h=\Delta^{1/4},\qquad \mres=n\sqrt\Delta,\\
 t_{\mathrm{cut}}=\frac{\log(1/\Delta)}{4\varkappa},\qquad
 \omega=e^{-t_{\mathrm{cut}}}=\Delta^{1/(4\varkappa)},\qquad
 L_\Delta=1+\log(1/\Delta).
 \end{gathered}
\end{equation}
We take $\Delta\le\Delta_0$ for a sufficiently small constant $\Delta_0>0$. Partition $[0,t_{\mathrm{cut}}]$ into equal intervals of widths between $h/2$ and $h$, and put $H_* =\{i:X_i>t_{\mathrm{cut}}\}$. A block is \emph{early} when its right endpoint is at most a fixed $\tau_0\in(0,1)$. If $b$ is the number of finite blocks and $b_{\mathrm E}$ the number of early blocks, then
\begin{equation}\label{eq:block-counts}
 b\le Ch^{-1}L_\Delta,\qquad c h^{-1}\le b_{\mathrm E}\le Ch^{-1}.
\end{equation}
Let $m_a=|H_a|=|Q_a|$ and $N_{ab}=|Q_a\cap H_b|$.

\begin{lemma}[Good clock environments]\label{lem:environment}
There is an event $\cG$ on which every finite row has $m_a\ge c\mres$, and, for every integer interval $J\subset[n]$ of length at least $c\mres$,
\begin{align}
 |J\cap H_b|&\ge ch|J| &&\text{for every early }b,\label{eq:early-cell}\\
 |J\cap H_*|&\le C(\omega|J|+\log n),
 & |H_*|&\le Cn\omega .\label{eq:tail-cell}
\end{align}
Uniformly in the parameters under consideration,
\begin{equation}\label{eq:environment-prob}
 \Pp(\cG^c)\le C\{h^{-1}L_\Delta e^{-c\mres}
     +n^2h^{-1}e^{-ch\mres}+n^{-2}\}=O_\varkappa(n^{-2}).
\end{equation}
Consequently, uniformly for $1\le t\le s$,
\begin{equation}\label{eq:tilted-environment}
 \Pp_t(\cG^c)\le C_\varkappa n^{-1+1/(4\pdual)}.
\end{equation}
\end{lemma}
\begin{proof}
For an interval of width between $h/2$ and $h$ inside $[0,t_{\mathrm{cut}}]$, each clock has probability at least $c h e^{-\varkappa t_{\mathrm{cut}}}=c\Delta^{1/2}$ of lying in it. Thus the expected row size is at least $c\mres$. For an early block each label has probability at least $ch$, and for the tail it has probability at most $\omega$. The relevant indicators are independent under the original measure. Multiplicative Chernoff bounds give the first two lower estimates, with failures $e^{-c\mres}$ and $e^{-ch|J|}$. For a Bernoulli sum of mean at most $\mu$, the elementary exponential bound
\[
 \Pp(Z\ge e\mu+x)\le e^{-x}
\]
follows by applying Markov's inequality to $e^Z$. Take $x=4\log n$ for the tail intervals and use at most $n^2$ intervals. A union bound and \eqref{eq:block-counts} give \eqref{eq:environment-prob}. The total-tail estimate follows in the same way, since $n\omega\gg\log n$.

Because $\Delta\ge1/n$, $\mres\ge\sqrt n$ and $h\mres=n\Delta^{3/4}\ge n^{1/4}$. The exponential terms are uniformly smaller than any fixed inverse power of $n$. Finally, H\"older with exponent $2$ and \eqref{eq:weight-norm} gives
$\Pp_t(\cG^c)\le\norm{W_{n,t}}_2\Pp(\cG^c)^{1/2}$, proving \eqref{eq:tilted-environment}.
\end{proof}

\subsection{Prefix depletion is distribution-free}
Choose a uniform label $I$, independent of the clocks under $\Pp_t$, and expose the first $k=r-\ell$ edges of its orbit. Write $v_j=\sigma_n^j(I)$ and define
\[
 V_a=\sum_{j<k}\one_{\{v_j\in H_a\}},\qquad
 U_{ab}=\sum_{j<k}\one_{\{v_j\in H_a,\ v_{j+1}\in H_b\}}.
\]
Let $\mathcal D$ be the event $V_a\le m_a/4$ for all finite $a$ and $U_{ab}\le N_{ab}/2$ for all finite $a$ and early $b$.
For every completed permutation and clock configuration, $v_j$ is uniform when only $I$ is random. Also
\[
 \#\{x\in H_a:\sigma_n(x)\in H_b\}=N_{ab}.
\]
Conditional Markov bounds followed by a union bound give, under any law of the clocks,
\begin{equation}\label{eq:depletion}
 \Pp_t(\mathcal D^c)\le C\Delta(b+bb_{\mathrm E})
 \le C\Delta^{1/2}L_\Delta.
\end{equation}
Zero-size cells contribute zero and need not be divided by. The probability that the prefix closes is at most $C_\varkappa\Delta$ by Corollary~\ref{cor:root-escape}. No tilted density norm is applied to \eqref{eq:depletion}.

On a good, nonclosing prefix, let $S_a^0\subset H_a$ be the unused sources and $T_a^0\subset Q_a$ the unused targets. Then
\[
 |S_a^0|=|T_a^0|=:m_a^0,\quad
 S^0=\bigcup_aS_a^0,\quad T^0=\bigcup_aT_a^0,\quad N_0=|S^0|=|T^0|,
\]
and
\begin{equation}\label{eq:residual-defects}
 T^0\setminus S^0=\{I\},\qquad S^0\setminus T^0=\{v_k\}.
\end{equation}
Put $N_{ab}^0=|T_a^0\cap S_b^0|$. On $\cG\cap\mathcal D$,
\begin{align}
 m_a^0&\ge\tfrac34m_a\ge c\mres &&(a\text{ finite}),\label{eq:residual-size}\\
 N_{ab}^0&\ge chm_a^0 &&(a\text{ finite},\ b\text{ early}),\label{eq:residual-early}\\
 N_{a*}^0/m_a^0&\le C(\omega+\log n/\mres),
 & m_*^0/N_0&\le C\omega .\label{eq:residual-tail}
\end{align}
Indeed, deleting prefix targets removes $U_{ab}$ points from the cell $Q_a\cap H_b$. Among surviving targets at most one further point has ceased to be a residual source, by \eqref{eq:residual-defects}. Thus $N_{ab}^0\ge N_{ab}-U_{ab}-1$. The remaining assertions follow from the environment estimates. Since $N_0\ge n-r$, its ratio to $n$ is uniformly bounded below.

\subsection{Reference reset and adaptive terminal exposure}
Let $\nu$ be uniform on the early block indices. For a finite row define
\[
 \widehat P^0(a,b)=\frac{N_{ab}^0}{m_a^0}
       +\frac{|T_a^0\setminus S^0|}{m_a^0}\nu_b.
\]
Use the same formula on a nonempty tail row, then replace that row by $\nu$ to obtain $\overline P^0$. Empty rows are also set to $\nu$. Define $p_a^0=m_a^0/N_0$. Equation~\eqref{eq:residual-early} gives a constant $\alpha>0$ such that
\begin{equation}\label{eq:minorization}
 \overline P^0(a,\cdot)\ge\alpha\nu(\cdot),\qquad
 \norm{\mu\overline P^0-\mu'\overline P^0}_{\TV}
 \le(1-\alpha)\norm{\mu-\mu'}_{\TV}.
\end{equation}
Exact column counting using \eqref{eq:residual-defects} gives
\[
 (p^0\widehat P^0)_b=
 \frac{|T^0\cap S_b^0|}{N_0}+\frac{\nu_b}{N_0},
\]
whereas $p_b^0$ differs from the first term by the row indicator of the tip divided by $N_0$. Therefore
\begin{equation}\label{eq:approx-invariance}
 \norm{p^0\overline P^0-p^0}_{\TV}\le N_0^{-1}+p_*^0
 \le C(\omega+N_0^{-1}).
\end{equation}
Writing $\mathbf e_a$ for the point mass at row $a$, iteration of \eqref{eq:minorization} and \eqref{eq:approx-invariance} yields
\begin{equation}\label{eq:reset-mixing}
 \norm{\mathbf e_a(\overline P^0)^j-p^0}_{\TV}
 \le(1-\alpha)^j+C(\omega+N_0^{-1}).
\end{equation}

During a compatible terminal history, the evolving residual sets satisfy
\[
 T^{(u)}\setminus S^{(u)}=\{I\},\qquad
 S^{(u)}\setminus T^{(u)}=\{\text{current tip}\}.
\]
This follows by deleting the current source and the chosen nonroot target at every step. In particular, a remaining target belonging to $S^0$ is still a remaining source. Its source-block index is therefore a legitimate next block state.

Assume the numerical conditions
\begin{equation}\label{eq:depth-conditions}
 1\le\ell\le r/3,\qquad \ell\le c_0\mres,\qquad \Delta\le\Delta_0,
\end{equation}
with $c_0$ sufficiently small. At most $\ell$ further targets are removed
from any frozen row, so its current size is at least $c\mres$. The minorization
is used only for the frozen reference kernel: no lower bound on every current
residual cell is needed after these terminal deletions. Lemma~\ref{lem:row-flatness} and comparison of the current and frozen uniform target sets give one-step total-variation error
\[
 C_\varkappa(h+\ell/\mres)
\]
relative to the completed reference transition. The coordinate upper bound \eqref{eq:row-upper} bounds a root hit by $C/\mres$ and a move to a tail source by $C(\omega+\log n/\mres)$. Complete a root hit artificially with $\nu$, but mark it as a coupling failure. Also mark any tail visit as a failure. A stepwise maximal coupling with the fixed reset chain has probability of failure at most
\begin{equation}\label{eq:terminal-failure}
 C_\varkappa\{\ell h+\ell^2/\mres+\ell\omega+\ell\log n/\mres\}.
\end{equation}
This comparison is uniform in every compatible history and does not suppose that the actual block process is Markov.

After $\ell-1$ block transitions, use \eqref{eq:reset-mixing}. If $\lambda_a^0=\Unif(T_a^0)$, the reference last-target law is exactly
\[
 \sum_a p_a^0\lambda_a^0=\Unif(T^0).
\]
The final row-flatness comparison has error $C(h+\ell/\mres)$, the reference tail row costs $C\omega$, and replacing $\Unif(T^0)$ by $U_n=\Unif([n])$ costs at most $r/n$.
The initial terminal tip is in the tail with unconditional probability at most $C\omega+\Pp_t(\cG^c)$: conditional on the full permutation, $\sigma_n^k(I)$ is uniform. Charging this event, \eqref{eq:depletion}, the prefix closure event, and \eqref{eq:tilted-environment} proves the following finite statement.

\begin{proposition}[One-root tilted endpoint interface]\label{prop:endpoint}
Under \eqref{eq:depth-conditions}, define
\begin{equation}\label{eq:endpoint-error}
 \begin{split}
 \mathcal E_{n,r,\ell}={}&(1-\alpha)^{\ell-1}
 +\ell h+\ell\omega+\frac{\ell^2+\ell\log n}{\mres}\\
 &+\omega+\Delta+\Delta^{1/2}L_\Delta+n^{-1/4}.
 \end{split}
\end{equation}
Then, for all sufficiently large $n$, uniformly in $1\le t\le s$,
\begin{equation}\label{eq:endpoint-TV}
 \E_t\norm{\mathcal L_t(\sigma_n^r(I)\mid\cH,I)-U_n}_{\TV}
 \le C_\varkappa\mathcal E_{n,r,\ell}.
\end{equation}
The constants are independent of $r,\ell,n,t$.
\end{proposition}

\begin{proof}
The conditional terminal comparison just established holds on each good, nonclosing prefix beginning its terminal segment in a finite row. Conditional on $\cH,I$, average it over the revealed prefix. Convexity of total variation permits charging the conditional probability of all excluded histories. Taking expectation, their total probability is bounded by \eqref{eq:tilted-environment}, \eqref{eq:depletion}, Corollary~\ref{cor:root-escape}, and the initial-tail bound. The term $n^{-1/4}$ dominates the polynomial environment error for the exponents in \eqref{eq:exponents}. The displayed coupling and mixing errors are exactly the other terms of \eqref{eq:endpoint-error}. If this bracket is bounded below by a fixed positive constant, the assertion also follows from the trivial total-variation bound. This accounts for parameter choices for which the right side is not small.
\end{proof}

\subsection{A truncation-uniform first-moment comparison}
For $T\ge1$, put
\begin{equation}\label{eq:truncations}
 Y_{n,r}^{(T)}=\frac1n\sum_{i:L_n(i)>r}(R_{i,\sigma_n^r(i)}\wedge T),
 \qquad
 V_{n,T}=\frac1{n^2}\sum_{i,j}(R_{ij}\wedge T),
 \quad\mu_{n,T}=\E V_{n,T}.
\end{equation}

\begin{lemma}[Truncated observable comparison]\label{lem:truncated}
Under \eqref{eq:depth-conditions}, for all sufficiently large $n$, uniformly for all $T\ge1$ and $1\le t\le s$,
\begin{align}
 \abs{\E_tY_{n,r}^{(T)}-\E_tV_{n,T}}
 &\le C_\varkappa T\mathcal E_{n,r,\ell},\label{eq:observable-compare}\\
 \E_t|V_{n,T}-\mu_{n,T}|
 &\le C_\varkappa Tn^{-1/2+1/(4\pdual)}.\label{eq:tilted-concentration}
\end{align}
\end{lemma}
\begin{proof}
For a finite clock block, replace a clock value by the block midpoint. If $i,j$ both lie in finite blocks, the resulting error in $\log R_{ij}$ is at most $2\varkappa h$. Since $u\mapsto e^u\wedge T$ is $T$-Lipschitz, the error in a factor is at most $2\varkappa Th$. The coarsened factor is a bounded $\cH$-measurable function of $i,j$. Apply \eqref{eq:endpoint-TV} to this function, with $i=I$, and replace the endpoint by a fresh independent uniform label $J$.

A tail root, a tail endpoint, or a tail fresh label has probability at most $C\omega+\Pp_t(\cG^c)$. For the endpoint this follows even without mixing: conditional on the full permutation and clocks, $\sigma_n^r(I)$ is uniform. Removing the indicator $L_n(I)>r$ costs $CT\Delta$ by \eqref{eq:cycle-crude}. Undoing the coarsening proves \eqref{eq:observable-compare}.

Under the original independent-clock law, replacing one clock changes $V_{n,T}$ by at most $2T/n$. The Doob martingale obtained by revealing the clocks therefore has squared increment sum at most $4T^2/n$, so
\[
 \Var(V_{n,T})\le4T^2/n.
\]
H\"older, \eqref{eq:weight-norm}, and $\pdual\ge2$ now give
\[
 \E_t|V_{n,T}-\mu_{n,T}|
 \le\norm{W_{n,t}}_2\sqrt{\Var(V_{n,T})}
 \le C_\varkappa Tn^{-1/2+1/(4\pdual)}.
\]
The bound is uniform in $T$; no tilted identity $\E_tR_{ij}=1$ has been used.
\end{proof}

\section{Summable real-tilted switching and optimal growth}\label{sec:real}

\begin{proposition}[Finite tilted mean-intensity bound]\label{prop:finite-mean}
Under \eqref{eq:depth-conditions}, for all sufficiently large $n$,
\begin{equation}\label{eq:finite-mean}
 \sup_{1\le t\le s}\abs{\E_t(S_{n,r}/n)-1}
 \le C_\varkappa\mathcal E_{n,r,\ell}^{\,\beta},
 \qquad \beta=\frac{p-1}{p}=\frac1\pdual>0.
\end{equation}
\end{proposition}
\begin{proof}
By Lemma~\ref{lem:selected},
\[
 \E_t(S_{n,r}/n-Y_{n,r}^{(T)})\le CT^{1-p},
 \qquad 0\le1-\mu_{n,T}\le CT^{1-p}.
\]
The second estimate uses the original independent-pair normalization $\E R_{ij}=1$, not a tilted normalization. Lemma~\ref{lem:truncated} therefore gives
\[
 \abs{\E_t(S_{n,r}/n)-1}
 \le C\{T\mathcal E_{n,r,\ell}+T^{1-p}\}.
\]
Since $\pdual\ge2$, the concentration term satisfies
\[
 Tn^{-1/2+1/(4\pdual)}\le Tn^{-1/4}\le T\mathcal E_{n,r,\ell},
\]
and is therefore absorbed in the displayed bound. If $\mathcal E\le1$, choose $T=\mathcal E^{-1/p}$. If $\mathcal E>1$, use the uniform first-moment bound on $S_{n,r}/n$. This proves \eqref{eq:finite-mean}.
\end{proof}

\begin{remark}
Only the absolute value of the \emph{mean} defect appears in \eqref{eq:finite-mean}. We neither assume nor need an $L^1(\Pp_t)$ concentration theorem for the switching intensity. This is why the one-root endpoint estimate suffices.
\end{remark}

\begin{proposition}[Summable real-tilted cycle error]\label{prop:summable}
There are constants $a,c,C>0$ depending only on $\varkappa$ such that, for all sufficiently large fixed $R$, sufficiently small fixed $\eps>0$, and all large $n$, uniformly for $R\le r\le\eps n$ and $1\le t\le s$,
\begin{equation}\label{eq:pointwise-tilted}
 \begin{split}
 |r\E_tC_{n,r}-t|\le C\bigg\{&e^{-cr}
 +(r/n)^a(1+\log(n/r))^C\\
 &+n^{-a}(1+\log n)^C+r/n\bigg\}.
 \end{split}
\end{equation}
Consequently,
\begin{equation}\label{eq:summable-tilted}
 \lim_{\eps\downarrow0}\lim_{R\to\infty}\limsup_{n\to\infty}
 \sup_{1\le t\le s}\sum_{R<r\le\eps n}
 \abs{\E_tC_{n,r}-t/r}=0.
\end{equation}
\end{proposition}
\begin{proof}
Choose $\ell=\lfloor\min\{r/3,\log(n/r)\}\rfloor$. For the stated range this is at least one, and \eqref{eq:depth-conditions} holds uniformly for large $n$, because $\mres=\sqrt{nr}\ge\sqrt n$ while $\ell\le\log n$. The reset term is bounded by $Ce^{-cr}+C(r/n)^c$. Every term involving $h,\omega,\Delta$ is a positive power of $r/n$ times a power of $1+\log(n/r)$. Also
\[
 \frac{\ell^2+\ell\log n}{\sqrt{nr}}
 \le Cn^{-1/2}(1+\log n)^2.
\]
Apply Proposition~\ref{prop:finite-mean}, reducing the positive exponents and enlarging $C$. The exact identity \eqref{eq:tilted-switch} gives
\[
 r\E_tC_{n,r}=t\frac n{n-r}\E_t(S_{n,r}/n),
\]
so the additional error is $O(r/n)$. This proves \eqref{eq:pointwise-tilted}.
After division by $r$, the profile-free scale term is bounded by
\[
 C\int_0^\eps x^{a-1}(1+\log(1/x))^C\dd x,
\]
which tends to zero with $\eps$. The exponential tail vanishes with $R$, the $n$-dependent term is $o(1)$ after harmonic summation, and the last term sums to $O(\eps)$. This proves \eqref{eq:summable-tilted}.
\end{proof}

\begin{corollary}[Optimal real growth]\label{cor:real-growth}
Uniformly over \eqref{eq:comparable},
\begin{equation}\label{eq:tilted-mean-total}
 \sup_{1\le t\le s}|\E_tK_n-tH_n|\le C_\varkappa,
 \qquad
 G_n(s)\le C_\varkappa n^{s-1}.
\end{equation}
\end{corollary}
\begin{proof}
Fix admissible $R$ and $\eps$. The sum of \eqref{eq:pointwise-tilted} divided by $r$ is bounded uniformly in $n,t$. For $r<R$, \eqref{eq:cycle-crude} bounds the total contribution by a constant depending only on the fixed $R,\varkappa$. The number of cycles longer than $\eps n$ is at most $1/\eps$, and $\sum_{r>\eps n}t/r$ is bounded. This proves the first assertion for large $n$; finitely many smaller $n$ are absorbed in the constant uniformly over their rate vectors. Since $G_n(1)=1$ and
\[
 \frac{d}{dt}\log G_n(t)=\frac{\E_tK_n}{t},
\]
integration gives $\log G_n(s)=(s-1)H_n+O_\varkappa(1)$, proving the second assertion.
\end{proof}

\section{An exact uniform component}\label{sec:uniform-core}

\begin{lemma}[Common-rate thinning]\label{lem:thinning}
Enlarge the clock probability space by independent auxiliary randomness and, conditionally on the clocks, take independent indicators with
\begin{equation}\label{eq:thinning}
 \Pp(B_i=1\mid X_i=x)=e^{-(\varkappa-\theta_i)x}.
\end{equation}
Set $S=\{i:B_i=1\}$, $N=|S|$, and let $D$ be the number of cycles of $\sigma_n$ avoiding $S$. Then
\begin{align}
 N&\ \overset d=\ \sum_{i=1}^n\operatorname{Bernoulli}(\theta_i/\varkappa)
     \quad\text{with independent summands},\label{eq:N-law}\\
 G_n(z)&=\E[z^D u_N(z)],\qquad
 u_m(z)=\prod_{j=1}^m\left(1+\frac{z-1}{j}\right),\quad u_0(z)=1.
 \label{eq:uniform-core}
\end{align}
There are constants $\xi_{\varkappa},c>0$ depending only on $\varkappa$ for which
\begin{equation}\label{eq:N-tail}
 \Pp(N<\xi_{\varkappa}n)\le e^{-cn}.
\end{equation}
\end{lemma}
\begin{proof}
The marked subdensity is
\[
 \Pp(B_i=1,X_i\in dx)=\theta_i e^{-\varkappa x}\dd x.
\]
Thus $\Pp(B_i=1)=\theta_i/\varkappa$ and, conditional on $B_i=1$, the clock has law $\Exp(\varkappa)$. The marked pairs are independent across labels. This proves \eqref{eq:N-law}, and the Chernoff lower-tail bound gives \eqref{eq:N-tail}, for example with $\xi_{\varkappa}=1/(2\varkappa)$.

Condition on $S$, all labeled clocks outside $S$, and the ordered unlabeled clocks inside $S$. The images of sources outside $S$ are fixed. The remaining map $\pi:S\to T$ to the unused ranks is a uniform bijection. Delete the outgoing edges of $S$ and contract the resulting fixed directed paths, including length-zero paths. This gives a bijection $\tau:T\to S$. Closed fixed components are exactly the $D$ cycles avoiding $S$. The permutation $g=\tau\circ\pi$ is uniform on $S$, and
\begin{equation}\label{eq:cycle-contraction}
 K_n=D+K(g).
\end{equation}
The number of cycles of a uniform permutation of $m$ objects has generating polynomial $u_m$. One proof inserts the largest element: it starts a new cycle in one of the $m$ equally likely insertion choices and otherwise joins an existing cycle, giving $u_m(z)=(1+(z-1)/m)u_{m-1}(z)$. Conditional expectation proves \eqref{eq:uniform-core}. When $N=0$, the identity means $D=K_n$. No independence of $D$ and $N$ is asserted or used.
\end{proof}

\begin{proposition}[A uniformly integrable residual count]\label{prop:residual-mgf}
For the $s>1$ in \eqref{eq:exponents},
\begin{equation}\label{eq:residual-exponential}
 \sup_{n,\,1\le\theta_i\le\varkappa}\E s^D<\infty.
\end{equation}
\end{proposition}
\begin{proof}
On $N\ge\xi_{\varkappa}n$, the real product formula gives $u_N(s)\ge c_{\varkappa,s}n^{s-1}$. Positivity in \eqref{eq:uniform-core} and Corollary~\ref{cor:real-growth} imply
\[
 \E[s^D;N\ge\xi_{\varkappa}n]\le C n^{1-s}G_n(s)\le C.
\]
On the complement, $D\le K_n$ and Cauchy--Schwarz gives
\[
 \E[s^D;N<\xi_{\varkappa}n]\le G_n(s^2)^{1/2}\Pp(N<\xi_{\varkappa}n)^{1/2}
 \le C n^{\varkappa(s^2-1)/2}e^{-cn/2}
\]
by Lemma~\ref{lem:crude-mgf}. This proves the assertion.
\end{proof}

\begin{proof}[Proof of Theorem~\ref{thm:uniform}]
Fix $0<\delta_\varkappa<\min\{(s-1)/4,1/4\}$. The expansion
\[
 \log\abs{1+(z-1)/j}=\Re(z-1)/j+O_{\delta_\varkappa}(j^{-2})
\]
gives, for all $m\ge0$ and $|z-1|\le\delta_\varkappa$,
\begin{equation}\label{eq:uniform-complex-bound}
 |u_m(z)|\le C_{\delta_\varkappa}(m+1)^{\Re z-1}.
\end{equation}
On $N\ge\xi_{\varkappa}n$, the modulus of $n^{1-z}z^D u_N(z)$ is at most $C_\varkappa s^D$. On $N<\xi_{\varkappa}n$, the factor $n^{1-\Re z}(N+1)^{\Re z-1}$ is at most $C n^{\delta_\varkappa}$; the exceptional expectation is therefore exponentially small by the preceding proof. Equations~\eqref{eq:uniform-core} and \eqref{eq:residual-exponential} prove \eqref{eq:normal-bound}.

The mean bound in \eqref{eq:uniform-cumulants} is the case $t=1$ of \eqref{eq:tilted-mean-total}. Put $F_n(z)=n^{1-z}G_n(z)$. A Cauchy derivative bound on a smaller disk, together with $F_n(1)=1$, yields a radius $r_\varkappa>0$ on which $|F_n(z)-1|\le1/2$ uniformly. Thus the branch of $\log F_n$ equal to zero at $1$ is uniformly bounded there. Cauchy's coefficient estimate implies
\[
 |\kappa_k^{\mathrm F}(K_n)|\le C_\varkappa k!r_\varkappa^{-k},\qquad k\ge2,
\]
since the normalization subtracts only the linear term $(z-1)\log n$. Absorb the factor $k$ into an exponential constant to obtain \eqref{eq:uniform-cumulants}.
\end{proof}

\section{The microscopic operator and the unweighted mean}\label{sec:operator}

We now assume \eqref{eq:profile} and use Inputs~\ref{input:micro}--\ref{input:root}. The operator assertions below also give a direct justification of the ordinary Fredholm determinant in the limiting residue; see Simon \cite[Chapter 3]{Simon} for trace-class determinants.

\begin{proposition}[Trace class, spectral gap, and trace correction]\label{prop:operator}
The kernel $\rho$ is doubly stochastic. The operator $\Kop$ is trace class and
\[
 \Tr\Kop^r=r\lambda_r\quad(r\ge1),\qquad
 \norm{\Kop-\Pi}_{2\to2}<1.
\]
The series and determinant formula in \eqref{eq:Af} hold, and
$|\lambda_r-1/r|\le C\vartheta^{r-1}/r$ for some $\vartheta<1$.
\end{proposition}
\begin{proof}
The substitution $y=\cA(t)$ gives $\dd y=\cB(t)\dd t$. Integrating \eqref{eq:rho} first in $y$ or first in $x$ shows that both marginals are one.
Let $m=\min f$, $M_f=\max f$, $c_f=(m+M_f)/2$, and $d_f=(M_f-m)/2$. Expand
\[
 \rho(x,y)=\sum_{k\ge0}a_k(x)b_k(y),\quad
 a_k(x)=\frac{f(x)(c_f-f(x))^k}{k!},\quad
 b_k(y)=\frac{\tau(y)^k e^{-c_f\tau(y)}}{\cB(\tau(y))}.
\]
Choose $0<\eta<m$. The set $E_\eta=\{x:f(x)\le m+\eta\}$ has positive measure, and
$\cB(t)\ge m|E_\eta|e^{-(m+\eta)t}$. Consequently
\[
 \norm{a_k}_2\le M_f d_f^k/k!,\qquad
 \norm{b_k}_2^2\le
 \frac{(2k)!}{m|E_\eta|(M_f-\eta)^{2k+1}}.
\]
Since $\sqrt{(2k)!}/k!\le2^k$,
\[
 \norm{a_k}_2\norm{b_k}_2
 \le C_\eta\left(\frac{M_f-m}{M_f-\eta}\right)^k.
\]
The ratio is less than one. The rank-one series therefore converges absolutely in trace norm and agrees with the kernel on compact target strips, hence with $\Kop$. Its diagonal is absolutely integrable term by term, giving $\Tr\Kop=\int\rho(x,x)\dd x$. For $r\ge2$, the usual Hilbert--Schmidt cyclic trace formula, or increasing compact-strip truncations and positivity, gives $\Tr\Kop^r=r\lambda_r$.

For the gap, put
\[
 k(y)=\frac{m e^{-M_f\tau(y)}}{\cB(\tau(y))},\qquad
 \alpha_0=\int_0^1k(y)\dd y=m/M_f.
\]
If $m=M_f$, $\Kop=\Pi$. Otherwise $\rho(x,y)\ge k(y)$, so
$\Kop=\one\otimes k+(1-\alpha_0)P$, where $P$ is a Markov kernel. For $h\perp\one$ and $c_h=\int kh$, the mean of $\Kop h$ is zero. Thus
\[
 \norm{\Kop h}_2^2\le\norm{\Kop h-c_h}_2^2
 \le(1-\alpha_0)\int_0^1(1-k(y))|h(y)|^2\dd y
 \le(1-\alpha_0)\norm h_2^2.
\]
The middle inequality is Jensen's inequality and the column-mass identity for $P$. Hence $\Top=\Kop-\Pi$ is trace class with $\norm\Top<1$. Double stochasticity gives $\Tr\Top^r=r\lambda_r-1$. Therefore
\[
 \sum_{r\ge1}\left|\lambda_r-\frac1r\right|
 \le\sum_{r\ge1}\frac{\norm\Top_1\norm\Top^{r-1}}r<\infty,
\]
and the trace-log expansion proves \eqref{eq:Af}. The determinant is real, continuous and nonzero along $I-u\Top$, $0\le u\le1$, so it is positive and the specified real logarithm is unambiguous.
\end{proof}

\begin{proposition}[Profile mean shift]\label{prop:mean-profile}
Under \eqref{eq:profile},
\begin{equation}\label{eq:mean-profile}
 \E K_n=H_n+A_f+o(1).
\end{equation}
\end{proposition}
\begin{proof}
Split $\E K_n-H_n$ at fixed $R$ and $\eps n$. For $r\le R$, Input~\ref{input:micro} gives convergence to $\sum_{r\le R}(\lambda_r-1/r)$. The middle contribution vanishes as $R\to\infty$ and $\eps\downarrow0$ by \eqref{eq:summable-tilted} at $t=1$.
For the macroscopic contribution, the exact size-biased identity is
\[
 \E\sum_{r>\eps n}C_{n,r}
 =\E\left[\frac{1}{L_n(I_n)/n}\one_{\{L_n(I_n)/n>\eps\}}\right].
\]
The integrand is bounded by $1/\eps$, with its sole interior discontinuity at $\eps$. Input~\ref{input:root} and the absence of an atom of the uniform law there imply convergence to $\int_\eps^1\dd x/x=\log(1/\eps)$. The harmonic sum over the same lengths has the same limit. Finally remove the two fixed cutoffs and use Proposition~\ref{prop:operator}. This argument uses only the unweighted uniform-root law, not a finite-deletion or tilted macroscopic theorem.
\end{proof}

\section{The residual Poisson cloud}\label{sec:marked}

Retain the thinning in Lemma~\ref{lem:thinning}. Put
\begin{equation}\label{eq:retention-mean}
 c_* =\frac1\varkappa\int_0^1f(x)\dd x\in[1/\varkappa,1].
\end{equation}
Then $N/n\to c_*$ in probability, by \eqref{eq:N-law} and \eqref{eq:profile}. Let $D_{n,r}$ count the $r$-cycles avoiding $S$, so that $D=\sum_r D_{n,r}$.

\subsection{Marked microscopic entries}
The subdensity corresponding to a label not selected by the thinning is
\begin{equation}\label{eq:unselected-density}
 \Pp(B_i=0,X_i\in dt)
 =\theta_i\left(e^{-\theta_i t}-e^{-\varkappa t}\right)\dd t.
\end{equation}
Define
\begin{align}
 \rho_{\mathrm{res}}(x,y)
 &=\rho(x,y)\left(1-e^{-(\varkappa-f(x))\tau(y)}\right)\label{eq:residual-kernel}\\
 &=\rho(x,y)-\frac{f(x)e^{-\varkappa\tau(y)}}{\cB(\tau(y))}.
 \nonumber
\end{align}
Write $\Qop$ for the operator with this kernel.

\begin{lemma}[Marked local entry]\label{lem:marked-local}
On every target-separated bulk region in Input~\ref{input:micro}, uniformly over all distinct source labels, for fixed $k$,
\begin{equation}\label{eq:marked-local}
 \Pp\{\sigma_n(i_a)=j_a,\ B_{i_a}=0,\ 1\le a\le k\}
 =n^{-k}\prod_{a=1}^k \rho_{\mathrm{res}}(i_a/n,j_a/n)+o(n^{-k})
\end{equation}
uniformly. The error is absolute, including when some limiting marked factor is zero.
\end{lemma}
\begin{proof}
Order the specified target ranks so that $j_1<\cdots<j_k$, and let $P=\{i_1,\ldots,i_k\}$. Conditional on $X_{i_a}=t_a$ with $t_1<\cdots<t_k$, the rank constraints are
\[
 \sum_{v\notin P}\one_{\{X_v\le t_a\}}=j_a-a,
 \qquad 1\le a\le k.
\]
Their joint probability, integrated against the product of the selected clock densities, is the unmarked entry probability. Adding the marks multiplies its integrand by
\[
 \prod_{a=1}^k\phi_{n,i_a}(t_a),\qquad
 \phi_{n,i}(t)=1-e^{-(\varkappa-\theta_{n,i})t}.
\]
Each factor lies in $[0,1]$ and is uniformly $\varkappa$-Lipschitz.

Let $t_{n,a}^*$ solve
\[
 \frac1n\sum_{v\notin P}(1-e^{-\theta_{n,v}t_{n,a}^*})=(j_a-a)/n.
\]
On separated bulk regions these saddles lie in a fixed compact interval and converge uniformly to $\tau(j_a/n)$; their derivatives are uniformly bounded below. This also follows directly from uniform convergence of the displayed empirical distribution functions. If any $|t_a-t_{n,a}^*|>n^{-1/4}$, a required marginal Bernoulli count differs from its mean by at least $c n^{3/4}$, or by a fixed multiple of $n$ outside a larger compact interval. Hoeffding's inequality \cite{Hoeffding} bounds that part of the integrated probability by $C_k e^{-c\sqrt n}$. On the complementary region, the product of marking factors differs from $\prod_a\phi_{n,i_a}(t_{n,a}^*)$ by $O_k(n^{-1/4})$.

The unmarked entry probability is $O(n^{-k})$ on the bulk region by Input~\ref{input:micro}. Therefore the marked probability equals
\[
 \left(\prod_a\phi_{n,i_a}(t_{n,a}^*)\right)
 \Pp\{\sigma_n(i_a)=j_a,\ 1\le a\le k\}
 +O(n^{-k-1/4})+O(e^{-c\sqrt n}).
\]
Now use \eqref{eq:local-input}, the uniform profile convergence, and the saddle convergence. No division by a marking factor is involved, proving the last assertion as well.
\end{proof}

\begin{proposition}[Fixed residual cycle counts]\label{prop:marked-cycles}
For each fixed $R$,
\begin{equation}\label{eq:fixed-residual}
 (D_{n,1},\ldots,D_{n,R})\Rightarrow(Z_1^-,\ldots,Z_R^-),
 \qquad Z_r^-\sim\Pois(\lambda_r^-)
\end{equation}
independently, with convergence of all fixed mixed factorial moments, where
\begin{equation}\label{eq:residual-traces}
 \lambda_r^-=\frac1r\int_{[0,1]^r}\rho_{\mathrm{res}}(x_1,x_2)\cdots \rho_{\mathrm{res}}(x_r,x_1)
 \dd x_1\cdots\dd x_r=\frac1r\Tr(\Qop^r).
\end{equation}
\end{proposition}
\begin{proof}
Expand a mixed factorial moment as a sum over ordered collections of distinct disjoint directed cycles. On the separated bulk configurations apply Lemma~\ref{lem:marked-local}. The Riemann sums factor over cycles; division by each cycle length accounts for cyclic choices of its starting vertex. On the complement, every marked summand is at most the corresponding unmarked summand, because all marking factors lie in $[0,1]$. For each fixed cycle encoding the target tuple determines the source tuple,
so the unmarked summand is bounded by the source supremum in
\eqref{eq:envelope-input}. That envelope removes the omitted tuples in the
same successive cutoff limits. This gives the product of the means in \eqref{eq:residual-traces}. Independent Poisson laws are moment-determinate, so the method of moments proves \eqref{eq:fixed-residual}; compare the general local-entry approach in \cite{Mukherjee}. The trace identity and finiteness are justified in the next lemma.
\end{proof}

\subsection{Subcriticality and a rank-one determinant identity}
Put
\begin{equation}\label{eq:rank-one}
 u(x)=f(x)/\varkappa,\qquad
 v(y)=\frac{\varkappa e^{-\varkappa\tau(y)}}{\cB(\tau(y))}.
\end{equation}
Then
\[
 \Qop=\Kop-u\otimes v,\quad
 \int u=c_*,\quad \int v=1,
\]
where $(u\otimes v)h=u\int bh$. Both $u$ and $v$ belong to $L^2$; for example $\cB(t)\ge e^{-\varkappa t}$ gives $\int v^2\le\varkappa$. The kernel $\rho_{\mathrm{res}}$ is nonnegative.

\begin{lemma}[Residual determinant]\label{lem:residual-det}
The operator $\Qop$ is trace class, $\norm{\Qop}_{2\to2}\le\sqrt{1-1/\varkappa}<1$ when $\varkappa>1$, and $\Qop=0$ when $\varkappa=1$. The means in \eqref{eq:residual-traces} are nonnegative and satisfy
\begin{equation}\label{eq:Lambda}
 \Lambda_f:=\sum_{r\ge1}\lambda_r^-
 =-\log\det(I-\Qop)=A_f-\log c_*<\infty.
\end{equation}
\end{lemma}
\begin{proof}
Trace class follows from Proposition~\ref{prop:operator} and the rank-one formula. Direct integration gives
\[
 \int \rho_{\mathrm{res}}(x,y)\dd y=1-f(x)/\varkappa\le1-1/\varkappa,
 \qquad
 0\le\int \rho_{\mathrm{res}}(x,y)\dd x\le1.
\]
Schur's test gives the asserted contraction. The diagonal trace formula follows by subtracting the rank-one kernel from the absolutely trace-convergent expansion of $\Kop$. For $r\ge2$, the Hilbert--Schmidt cyclic trace formula applies as before. In particular all the cyclic integrals are finite and nonnegative. The convergent trace-log expansion gives $\sum_r\lambda_r^-=-\log\det(I-\Qop)$.

For the constant, decompose $L^2[0,1]=\mathbb C\one\oplus L^2_0$ and write
\[
 u=c_*\one+u_0,\qquad v=\one+v_0,
 \qquad \Top_0=\Kop|_{L^2_0}.
\]
Double stochasticity makes $\Kop$ block diagonal with entries $1,\Top_0$. Relative to this decomposition,
\begin{equation}\label{eq:block-det}
 I-\Qop=
 \begin{pmatrix}
 c_*&c_*v_0^*\\
 u_0&I-\Top_0+u_0v_0^*
 \end{pmatrix},
\end{equation}
where $v_0^*$ denotes the functional $h\mapsto\int v_0h$. Multiply on the left by the identity plus the rank-one nilpotent block with lower-left entry $-u_0/c_*$. Its Fredholm determinant is one. The result is upper triangular with diagonal blocks $c_*$ and $I-\Top_0$. Multiplicativity of trace-class determinants yields
\begin{equation}\label{eq:det-factorization}
 \det(I-\Qop)=c_*\det(I-\Top_0)
             =c_*\det\bigl(I-(\Kop-\Pi)\bigr).
\end{equation}
Both determinants are positive by their real contraction paths, so taking real logarithms and using \eqref{eq:Af} proves \eqref{eq:Lambda}.
\end{proof}

\subsection{No residual mass escapes to growing cycle lengths}
\begin{proposition}[The full residual count]\label{prop:residual-Poisson}
Under \eqref{eq:profile},
\begin{equation}\label{eq:D-Poisson}
 D\Rightarrow\Pois(\Lambda_f).
\end{equation}
For every $0<s_0<s$, the generating functions converge uniformly on $|z|\le s_0$:
\begin{equation}\label{eq:D-pgf}
 \E z^D\longrightarrow e^{\Lambda_f(z-1)}.
\end{equation}
\end{proposition}
\begin{proof}
The conditional cycle decomposition gives
\begin{equation}\label{eq:D-mean-identity}
 \E D=\E K_n-\E H_N,\qquad H_0=0.
\end{equation}
By \eqref{eq:N-law}, $N/n\to c_*$ in probability. On $N\ge\xi_{\varkappa}n$,
$H_N-H_n=\log(N/n)+O(1/n)$ uniformly, and the logarithm is bounded. The complementary contribution is at most $H_n\Pp(N<\xi_{\varkappa}n)=o(1)$. Hence
\[
 \E H_N=H_n+\log c_*+o(1).
\]
Proposition~\ref{prop:mean-profile} and Lemma~\ref{lem:residual-det} imply
\begin{equation}\label{eq:D-mean-limit}
 \E D\longrightarrow A_f-\log c_*=\Lambda_f.
\end{equation}
For each fixed $R$, Proposition~\ref{prop:marked-cycles} gives
\[
 \E\sum_{r\le R}D_{n,r}\longrightarrow\sum_{r\le R}\lambda_r^-.
\]
Subtracting from \eqref{eq:D-mean-limit} and using nonnegativity,
\[
 \limsup_{n\to\infty}\Pp\left(D\ne\sum_{r\le R}D_{n,r}\right)
 \le\Lambda_f-\sum_{r\le R}\lambda_r^-\longrightarrow0.
\]
The fixed residual sum tends to $\Pois(\sum_{r\le R}\lambda_r^-)$, which tends to $\Pois(\Lambda_f)$ as $R\to\infty$. This proves \eqref{eq:D-Poisson}. Finally, \eqref{eq:residual-exponential} bounds the tails of the generating series on $|z|\le s_0<s$ by $C(s_0/s)^J$ beyond degree $J$. The finitely many probabilities of values at most $J$ converge by \eqref{eq:D-Poisson}. Letting $J\to\infty$ proves \eqref{eq:D-pgf}. This is an exponential-uniform-integrability argument, not a passage from total variation to an unbounded test function.
\end{proof}

\begin{proof}[Proof of Theorem~\ref{thm:profile}]
Proposition~\ref{prop:operator} proves the operator and determinant assertions. The radius already chosen in the proof of Theorem~\ref{thm:uniform} satisfies $s_0=1+\delta_\varkappa<s$. On $N\ge\xi_{\varkappa}n$, uniformly for $|z-1|\le\delta_\varkappa$, the uniform-permutation product has the gamma-ratio asymptotic \cite[Section 5.11(iii)]{DLMF}
\begin{equation}\label{eq:gamma-ratio}
 u_N(z)=\frac{\Gamma(N+z)}{\Gamma(z)\Gamma(N+1)}
       =\frac{N^{z-1}}{\Gamma(z)}\{1+O_{\delta_\varkappa}(N^{-1})\}.
\end{equation}
The factor $(N/n)^{z-1}$ converges uniformly in $z$ in probability to $c_*^{z-1}$ and is uniformly bounded on this event. To justify multiplication by $z^D$, choose $a_0>1$ so that $s_0^{a_0}<s$. Proposition~\ref{prop:residual-mgf} bounds $\E s_0^{a_0D}$ uniformly; H\"older and the bounded convergence in probability of $N/n$ then give
\[
 \sup_{|z-1|\le\delta_\varkappa}
 \abs{\E\left[z^D\{(N/n)^{z-1}-c_*^{z-1}\};N\ge\xi_{\varkappa}n\right]}
 \longrightarrow0.
\]
No independence is needed. The gamma-ratio error is $O(1/n)$ after expectation by the same exponential bound. On $N<\xi_{\varkappa}n$, \eqref{eq:uniform-complex-bound}, Cauchy--Schwarz, \eqref{eq:crude-mgf}, and \eqref{eq:N-tail} make the normalized contribution exponentially small, uniformly in $z$.
Using \eqref{eq:uniform-core} and \eqref{eq:D-pgf}, we obtain
\[
 n^{1-z}G_n(z)\longrightarrow
 \frac{c_*^{z-1}e^{\Lambda_f(z-1)}}{\Gamma(z)}
 =\frac{e^{A_f(z-1)}}{\Gamma(z)}.
\]
The convergence is uniform on the displayed closed disk and hence locally uniform on its interior. Equation~\eqref{eq:mean-profile}, together with $H_n-\log n\to\gamma$, gives \eqref{eq:exact-mean-residue}.
\end{proof}

\section{Consequences and scope}\label{sec:consequences}

\begin{corollary}[Universal higher factorial cumulants]\label{cor:limiting-cumulants}
Under \eqref{eq:profile}, for every fixed integer $k\ge2$,
\begin{equation}\label{eq:limiting-cumulants}
 \kappa_k^{\mathrm F}(K_n)\longrightarrow
 (-1)^{k+1}(k-1)!\zeta(k).
\end{equation}
In particular,
\[
 \Var K_n=\log n+\gamma+A_f-\frac{\pi^2}{6}+o(1).
\]
\end{corollary}
\begin{proof}
The quotient $G_n(z)/u_n(z)$ tends locally uniformly to $e^{A_f(z-1)}$ by Theorem~\ref{thm:profile} and \eqref{eq:gamma-ratio}. On a smaller disk all these functions admit the logarithm normalized at $1$, so every fixed derivative of $\log G_n-\log u_n$ of order at least two tends to zero. The exact product defining $u_n$ gives
\[
 \left.\frac{d^k}{dw^k}\log u_n(1+w)\right|_{w=0}
 =(-1)^{k+1}(k-1)!\sum_{j=1}^n j^{-k}.
\]
Let $n\to\infty$. The variance identity is $\Var K_n=\E K_n+\varkappa_2^{\mathrm F}(K_n)$.
\end{proof}

\begin{remark}[The choice of the extraction rate]
The upper bound $\varkappa$ used for thinning need not be the smallest available one. Different valid choices change $c_*$ and the residual Poisson parameter $\Lambda_f$ separately. Equation~\eqref{eq:Lambda} shows that $\Lambda_f+\log c_*=A_f$ is unchanged. In the constant-profile case $f\equiv a\le\varkappa$, $\rho\equiv1$, $A_f=0$, $c_*=a/\varkappa$, and $\Qop$ is rank one with eigenvalue $1-c_*$. Thus $\lambda_r^-=(1-c_*)^r/r$ and $\Lambda_f=-\log c_*$, recovering the uniform-permutation residue $1/\Gamma(z)$.
\end{remark}

\begin{remark}[What is and is not required]
The uniform analytic bound and the all-order estimate \eqref{eq:uniform-cumulants} hold for arbitrary comparable-rate arrays. The identification of the limiting residue in Theorem~\ref{thm:profile} uses regular-profile microscopic information and the unweighted macroscopic root law. No convergence rate uniform over all continuous profiles or over perturbation arrays without a common modulus is asserted. No claim of an identified exact-mean residue for every arbitrary comparable array is made here. The proof uses real fugacities only on the one-sided interval $[1,s]$; it obtains complex convergence from the exact uniform component rather than from positivity under a complex change of measure.
\end{remark}

\section*{Dependency summary}
The new uniform argument consists of the exact switch, the tilted row-insertion lemma, the finite one-root endpoint comparison, the optimal real generating-function bound, and the exact common-rate component. Their proofs are given above; the exact switch is also cited to \cite{LongMeso} for provenance rather than imported without proof. For the profile identification, the only imported asymptotic statements are Input~\ref{input:micro} and Input~\ref{input:root}. The marked microscopic extension, trace-class calculation, mean-shift passage, rank-one determinant identity, and residual-tail removal are proved in this manuscript. In particular, the conclusion is not derived from the higher-cumulant assertion that it subsequently implies.

\clearpage
\begin{thebibliography}{99}
\footnotesize\raggedright

\bibitem{BCD}
J.~Borga, S.~Chatterjee, and P.~Diaconis,
\emph{Permuton and local limits for the Luce model},
arXiv:2509.07729v3, 29 September 2026; accepted for publication in
Ann. Appl. Probab. See Section~5.2 for the cycle-structure questions.
\url{https://arxiv.org/abs/2509.07729v3}.

\bibitem{CDK}
S.~Chatterjee, P.~Diaconis, and G.~B.~Kim,
\emph{Enumerative theory for the Tsetlin library},
J. Algebra \textbf{655} (2024), 139--162.
\url{https://doi.org/10.1016/j.jalgebra.2023.08.009}.

\bibitem{DiaconisSeminar}
P.~Diaconis,
\emph{On sampling from an urn, without replacement},
Stanford Combinatorics Seminar, 4 December 2025, seminar abstract.
\url{https://mathematics.stanford.edu/events/sampling-urn-without-replacement}.

\bibitem{DLMF}
R.~A.~Askey and R.~Roy,
\emph{Gamma function}, Chapter 5 of the NIST Digital Library of Mathematical Functions,
Section 5.11(iii), equations 5.11.12--5.11.13.
\url{https://dlmf.nist.gov/5.11.iii}.

\bibitem{Hoeffding}
W.~Hoeffding,
\emph{Probability inequalities for sums of bounded random variables},
J. Amer. Statist. Assoc. \textbf{58} (1963), no.~301, 13--30.
\url{https://doi.org/10.1080/01621459.1963.10500830}.

\bibitem{KN}
E.~Kowalski and A.~Nikeghbali,
\emph{Mod-Poisson convergence in probability and number theory},
Int. Math. Res. Not. IMRN \textbf{2010} (2010), no.~18, 3549--3587.
\url{https://doi.org/10.1093/imrn/rnq019}.

\bibitem{LongMeso}
C.~D.~Long,
\emph{Mesoscopic cycles of comparable-rate Luce permutations: a switching proof of Conjecture 13.1},
research manuscript, 2026.
Repository \texttt{long-mathematics/luce-tsetlin-permutations},
revision \texttt{01f99e121c4e}, file
\path{papers/mesoscopic-switching/mesoscopic-luce-cycles.tex}.
\href{https://github.com/long-mathematics/luce-tsetlin-permutations/tree/01f99e121c4ea497e3c9e9bcfa83d027c5833742/papers/mesoscopic-switching}{Pinned repository source}.

\bibitem{LongMicro}
C.~D.~Long,
\emph{Microscopic cycles and sparse entries for regular Luce profiles},
research manuscript, 2026.
Repository \texttt{long-mathematics/luce-tsetlin-permutations},
revision \texttt{01f99e121c4e}, file
\path{papers/regular-profile-cycles/regular-profile-cycles.tex}.
\href{https://github.com/long-mathematics/luce-tsetlin-permutations/tree/01f99e121c4ea497e3c9e9bcfa83d027c5833742/papers/regular-profile-cycles}{Pinned repository source}.

\bibitem{LongMacro}
C.~D.~Long,
\emph{Poisson--Dirichlet universality for Luce permutations},
research manuscript, 11 October 2026.
Repository \texttt{long-mathematics/luce-tsetlin-permutations},
revision \texttt{01f99e121c4e}, file
\path{papers/macroscopic-pd/poisson-dirichlet-luce.tex}.
\href{https://github.com/long-mathematics/luce-tsetlin-permutations/tree/01f99e121c4ea497e3c9e9bcfa83d027c5833742/papers/macroscopic-pd}{Pinned repository source}.

\bibitem{LongTotal}
C.~D.~Long,
\emph{Poisson approximation and lattice Edgeworth expansions for total Luce cycles},
research manuscript, 11 October 2026.
Repository \texttt{long-mathematics/luce-tsetlin-permutations},
revision \texttt{01f99e121c4e}, file
\path{papers/total-cycle-edgeworth/total-cycle-edgeworth.tex}.
\href{https://github.com/long-mathematics/luce-tsetlin-permutations/tree/01f99e121c4ea497e3c9e9bcfa83d027c5833742/papers/total-cycle-edgeworth}{Pinned repository source}.

\bibitem{LongMean}
C.~D.~Long,
\emph{Profile-dependent mean shifts and universal factorial cumulants for Luce permutations},
Unpublished manuscript, 13 October 2026, revised 20 October 2026.
Cited for context; no result from this manuscript is an imported input here.

\bibitem{Luce}
R.~D.~Luce,
\emph{Individual Choice Behavior: A Theoretical Analysis},
John Wiley \& Sons, New York, 1959.

\bibitem{Mukherjee}
S.~Mukherjee,
\emph{Fixed points and cycle structure of random permutations},
Electron. J. Probab. \textbf{21} (2016), paper no.~40, 1--18.
\url{https://doi.org/10.1214/16-EJP4622}.

\bibitem{NZ}
A.~Nikeghbali and D.~Zeindler,
\emph{The generalized weighted probability measure on the symmetric group and the asymptotic behavior of the cycles},
Ann. Inst. Henri Poincar\'e Probab. Stat. \textbf{49} (2013), no.~4, 961--981.
\url{https://doi.org/10.1214/12-AIHP484}.

\bibitem{Pitman}
J.~Pitman,
\emph{Combinatorial Stochastic Processes},
Lecture Notes in Mathematics, vol.~1875, Springer, Berlin, 2006.
\url{https://doi.org/10.1007/b11601500}.

\bibitem{Simon}
B.~Simon,
\emph{Trace Ideals and Their Applications}, 2nd ed.,
Mathematical Surveys and Monographs, vol.~120,
American Mathematical Society, Providence, RI, 2005.
\url{https://doi.org/10.1090/surv/120}.

\bibitem{vBE}
B.~von Bahr and C.-G.~Esseen,
\emph{Inequalities for the $r$th absolute moment of a sum of random variables, $1\le r\le2$},
Ann. Math. Statist. \textbf{36} (1965), no.~1, 299--303.
\url{https://doi.org/10.1214/aoms/1177700291}.

\end{thebibliography}
\end{document}
